\section{Kết luận}

Nghiên cứu đã làm rõ tiềm năng và phương thức ứng dụng trí tuệ nhân tạo tạo sinh trong việc hỗ trợ giảng viên Đại học Bách khoa Hà Nội biên soạn học liệu và thiết kế hoạt động dạy học. Dựa trên kết quả khảo sát thực tế nhu cầu cùng những áp lực về thời gian soạn bài của giảng viên, bài luận đã đề xuất khung nguyên tắc thiết kế và phát triển công cụ thử nghiệm tích hợp các mô hình ngôn ngữ lớn để tự động hóa khâu tóm tắt tài liệu, tạo dàn ý bài giảng và sinh câu hỏi đánh giá.

Nghiên cứu khẳng định rằng việc xây dựng các công cụ hỗ trợ dạy học dựa trên trí tuệ nhân tạo phải tuân thủ nghiêm ngặt các nguyên lý sư phạm chuẩn bao gồm Thang đo nhận thức Bloom cải tiến và Thuyết kiến tạo, thay vì chỉ sử dụng công nghệ một cách tự phát. Đồng thời, mô hình hoạt động phù hợp nhất là mô hình có sự giám sát chuyên môn của con người, trong đó công nghệ đóng vai trò trợ lý xây dựng bản thảo thô và giảng viên giữ vai trò quyết định, kiểm duyệt tính chính xác của học liệu đầu ra. Hướng tiếp cận này giúp bảo đảm sự cân bằng giữa việc tối ưu hóa thời gian chuẩn bị bài giảng và việc duy trì liêm chính học thuật cùng tính chính xác của kiến thức chuyên ngành.

Về định hướng phát triển sản phẩm công nghệ trong tương lai, công cụ cần tập trung vào việc tích hợp trực tiếp vào các hệ thống quản lý học tập hiện có của nhà trường như Microsoft Teams hoặc Moodle để tạo sự thuận tiện tối đa cho người sử dụng. Bên cạnh đó, hệ thống cũng cần được nghiên cứu mở rộng tính năng để hỗ trợ xử lý học liệu đa phương tiện, giúp tự động hóa việc thiết kế sơ đồ trực quan và cấu trúc bản trình chiếu bài giảng, đáp ứng tốt hơn nhu cầu giảng dạy kỹ thuật chuyên sâu.
